\documentclass[10pt,a4paper]{amsart}
\usepackage{amsmath,amssymb,amsthm,graphicx}
% Retained preamble for the Stacks Project chapter snapshot in this directory.
% The chapter's own preamble.tex is a copy of the shared Stacks Project preamble; the block
% below is sliced MECHANICALLY, line for line, from it. Nothing was added, removed or
% redefined; the wrapper class, the xr-hyper/hyperref wiring and the externaldocument list
% are omitted because the assembler supplies the document class.
%
% lines taken from preamble.tex: 3, 5, 9, 11, 12, 156, 157, 158, ... (51 lines total)
%
% package coverage restored from the chapter's own preamble (2026-10-13): xy, xy-2cell

\usepackage[all]{xy}
\xyoption{2cell}
\UseAllTwocells

\newenvironment{reference}{\comment}{\endcomment}
\newenvironment{slogan}{\comment}{\endcomment}
\newenvironment{history}{\comment}{\endcomment}
\theoremstyle{plain}
\newtheorem{theorem}[subsection]{Theorem}
\newtheorem{proposition}[subsection]{Proposition}
\newtheorem{lemma}[subsection]{Lemma}
\theoremstyle{definition}
\newtheorem{definition}[subsection]{Definition}
\newtheorem{example}[subsection]{Example}
\newtheorem{exercise}[subsection]{Exercise}
\newtheorem{situation}[subsection]{Situation}
\theoremstyle{remark}
\newtheorem{remark}[subsection]{Remark}
\newtheorem{remarks}[subsection]{Remarks}
\numberwithin{equation}{subsection}
\def\lim{\mathop{\mathrm{lim}}\nolimits}
\def\colim{\mathop{\mathrm{colim}}\nolimits}
\def\Spec{\mathop{\mathrm{Spec}}}
\def\Hom{\mathop{\mathrm{Hom}}\nolimits}
\def\Ext{\mathop{\mathrm{Ext}}\nolimits}
\def\SheafHom{\mathop{\mathcal{H}\!\mathit{om}}\nolimits}
\def\SheafExt{\mathop{\mathcal{E}\!\mathit{xt}}\nolimits}
\def\Sch{\mathit{Sch}}
\def\Mor{\mathop{\mathrm{Mor}}\nolimits}
\def\Ob{\mathop{\mathrm{Ob}}\nolimits}
\def\Sh{\mathop{\mathit{Sh}}\nolimits}
\def\NL{\mathop{N\!L}\nolimits}
\def\CH{\mathop{\mathrm{CH}}\nolimits}
\def\proetale{{pro\text{-}\acute{e}tale}}
\def\etale{{\acute{e}tale}}
\def\QCoh{\mathit{QCoh}}
\def\Ker{\mathop{\mathrm{Ker}}}
\def\Im{\mathop{\mathrm{Im}}}
\def\Coker{\mathop{\mathrm{Coker}}}
\def\Coim{\mathop{\mathrm{Coim}}}
\DeclareMathSymbol{\boxtimes}{\mathbin}{AMSa}{"02}
\def\QCohstack{\mathcal{QC}\!\mathit{oh}}
\def\Cohstack{\mathcal{C}\!\mathit{oh}}
\def\Spacesstack{\mathcal{S}\!\mathit{paces}}
\def\Quotfunctor{\mathrm{Quot}}
\def\Hilbfunctor{\mathrm{Hilb}}
\def\Curvesstack{\mathcal{C}\!\mathit{urves}}
\def\Polarizedstack{\mathcal{P}\!\mathit{olarized}}
\def\Complexesstack{\mathcal{C}\!\mathit{omplexes}}
\def\Pic{\mathop{\mathrm{Pic}}\nolimits}
\def\Picardstack{\mathcal{P}\!\mathit{ic}}
\def\Picardfunctor{\mathrm{Pic}}
\def\Deformationcategory{\mathcal{D}\!\mathit{ef}}
\usepackage{cleveref}
\crefname{originallemma}{original lemma}{original lemmas}
\Crefname{originallemma}{Original lemma}{Original lemmas}
\makeatletter
\newlabel{algebra-lemma-characterize-flat}{{00M5}{source}{}{sourcepointer-1}{}}
\newlabel{algebra-lemma-characterize-flat@cref}{{[originallemma][0][]00M5}{[1][0][]source}{}{}{}}
\makeatother
\crefname{originallemma}{original lemma}{original lemmas}
\Crefname{originallemma}{Original lemma}{Original lemmas}
\makeatletter
\newlabel{algebra-lemma-what-does-it-mean}{{051C}{source}{}{sourcepointer-2}{}}
\newlabel{algebra-lemma-what-does-it-mean@cref}{{[originallemma][0][]051C}{[1][0][]source}{}{}{}}
\makeatother
\crefname{originallemma}{original lemma}{original lemmas}
\Crefname{originallemma}{Original lemma}{Original lemmas}
\makeatletter
\newlabel{algebra-remark-Tor-ring-mod-ideal}{{00M6}{source}{}{sourcepointer-3}{}}
\newlabel{algebra-remark-Tor-ring-mod-ideal@cref}{{[originallemma][0][]00M6}{[1][0][]source}{}{}{}}
\makeatother
\crefname{originallemma}{original lemma}{original lemmas}
\Crefname{originallemma}{Original lemma}{Original lemmas}
\makeatletter
\newlabel{definition-strict-morphism-thickenings}{{08L1}{source}{}{sourcepointer-4}{}}
\newlabel{definition-strict-morphism-thickenings@cref}{{[originallemma][0][]08L1}{[1][0][]source}{}{}{}}
\makeatother
\crefname{originallemma}{original lemma}{original lemmas}
\Crefname{originallemma}{Original lemma}{Original lemmas}
\makeatletter
\newlabel{modules-lemma-stalk-tensor-product}{{01CB}{source}{}{sourcepointer-5}{}}
\newlabel{modules-lemma-stalk-tensor-product@cref}{{[originallemma][0][]01CB}{[1][0][]source}{}{}{}}
\makeatother
\crefname{originallemma}{original lemma}{original lemmas}
\Crefname{originallemma}{Original lemma}{Original lemmas}
\makeatletter
\newlabel{sheaves-lemma-stalk-pullback-modules}{{0098}{source}{}{sourcepointer-6}{}}
\newlabel{sheaves-lemma-stalk-pullback-modules@cref}{{[originallemma][0][]0098}{[1][0][]source}{}{}{}}
\makeatother
\crefname{originallemma}{original lemma}{original lemmas}
\Crefname{originallemma}{Original lemma}{Original lemmas}
\makeatletter
\newlabel{situation-morphism-thickenings}{{08KZ}{source}{}{sourcepointer-7}{}}
\newlabel{situation-morphism-thickenings@cref}{{[originallemma][0][]08KZ}{[1][0][]source}{}{}{}}
\makeatother
\title{Stacks, Deformation Theory: lemma-deform-module}
\author{The Stacks Project Authors}
\date{stacks-project:defos.tex:local-snapshot-20261012; GNU FDL 1.2 or later}
\begin{document}
\maketitle
\noindent\textit{Editorial note.} Statement span is lines 1535-1559 of sources/stacks-defos/pinned-original-source.tex; proof span is lines 1561-1608, the complete author proof (46 body lines).\par\medskip
\noindent\textit{Editorial note.} Navigation stubs: 7 cited-but-not-included labels are supplied as non-load-bearing pointers with their pinned tags where the snapshot tag index has them.\par\medskip
\noindent\textit{Editorial note.} extra\_packages is empty on purpose: the retained preamble already loads xy with options and a second declaration clashes in the assembler wrapper.\par\medskip
\noindent\textit{Editorial note.} Text is retained verbatim, including the author's own typos and placeholders.\par\medskip
\noindent\textit{Editorial note.} Queue index 17 of queue\_defos.json (pb=1826, stmt 1535-1559, proof 1561-1608).\par\medskip
\noindent\textit{Allowed standard prerequisites (editorial).}\begin{itemize}
\item Explicit prerequisite boundary. The cited inputs are prerequisites; the argument in the retained span is not replaced by any cited result.
\end{itemize}\medskip
\ifdefined\hypertarget\hypertarget{sourcepointer-1}{}\fi
\noindent\textit{Source reference.} Original lemma 00M5: sources/stacks-defos/pinned-original-source.tex, source lines 1608--1608. Cited by the retained statement or proof but not part of the collected spans; supplied only as a navigation pointer. The argument does not depend on its content.\par\medskip
\ifdefined\hypertarget\hypertarget{sourcepointer-2}{}\fi
\noindent\textit{Source reference.} Original lemma 051C: sources/stacks-defos/pinned-original-source.tex, source lines 1608--1608. Cited by the retained statement or proof but not part of the collected spans; supplied only as a navigation pointer. The argument does not depend on its content.\par\medskip
\ifdefined\hypertarget\hypertarget{sourcepointer-3}{}\fi
\noindent\textit{Source reference.} Original lemma 00M6: sources/stacks-defos/pinned-original-source.tex, source lines 1608--1608. Cited by the retained statement or proof but not part of the collected spans; supplied only as a navigation pointer. The argument does not depend on its content.\par\medskip
\ifdefined\hypertarget\hypertarget{sourcepointer-4}{}\fi
\noindent\textit{Source reference.} Original lemma 08L1: sources/stacks-defos/pinned-original-source.tex, source lines 766--770. Cited by the retained statement or proof but not part of the collected spans; supplied only as a navigation pointer. The argument does not depend on its content.\par\medskip
\ifdefined\hypertarget\hypertarget{sourcepointer-5}{}\fi
\noindent\textit{Source reference.} Original lemma 01CB: sources/stacks-defos/pinned-original-source.tex, source lines 1608--1608. Cited by the retained statement or proof but not part of the collected spans; supplied only as a navigation pointer. The argument does not depend on its content.\par\medskip
\ifdefined\hypertarget\hypertarget{sourcepointer-6}{}\fi
\noindent\textit{Source reference.} Original lemma 0098: sources/stacks-defos/pinned-original-source.tex, source lines 1608--1608. Cited by the retained statement or proof but not part of the collected spans; supplied only as a navigation pointer. The argument does not depend on its content.\par\medskip
\ifdefined\hypertarget\hypertarget{sourcepointer-7}{}\fi
\noindent\textit{Source reference.} Original lemma 08KZ: sources/stacks-defos/pinned-original-source.tex, source lines 737--747. Cited by the retained statement or proof but not part of the collected spans; supplied only as a navigation pointer. The argument does not depend on its content.\par\medskip
\begin{lemma}
\label{lemma-deform-module}
Let $(f, f')$ be a morphism of first order thickenings of ringed spaces
as in Situation \ref{situation-morphism-thickenings}.
Let $\mathcal{F}'$ be an $\mathcal{O}_{X'}$-module
and set $\mathcal{F} = i^*\mathcal{F}'$.
Assume that $\mathcal{F}$ is flat over $S$
and that $(f, f')$ is a strict morphism of thickenings
(Definition \ref{definition-strict-morphism-thickenings}).
Then the following are equivalent
\begin{enumerate}
\item $\mathcal{F}'$ is flat over $S'$, and
\item the canonical map
$f^*\mathcal{J} \otimes_{\mathcal{O}_X} \mathcal{F} \to
\mathcal{I}\mathcal{F}'$
is an isomorphism.
\end{enumerate}
Moreover, in this case the maps
$$
f^*\mathcal{J} \otimes_{\mathcal{O}_X} \mathcal{F} \to
\mathcal{I} \otimes_{\mathcal{O}_X} \mathcal{F} \to
\mathcal{I}\mathcal{F}'
$$
are isomorphisms.
\end{lemma}

\begin{proof}
The map $f^*\mathcal{J} \to \mathcal{I}$ is surjective
as $(f, f')$ is a strict morphism of thickenings.
Hence the final statement is a consequence of (2).

\medskip\noindent
Proof of the equivalence of (1) and (2). We may check these conditions
at stalks. Let $x \in X \subset X'$
be a point with image $s = f(x) \in S \subset S'$.
Set $A' = \mathcal{O}_{S', s}$, $B' = \mathcal{O}_{X', x}$,
$A = \mathcal{O}_{S, s}$, and $B = \mathcal{O}_{X, x}$.
Then $A = A'/J$ and $B = B'/I$ for some square zero ideals.
Since $(f, f')$ is a strict morphism of thickenings we have $I = JB'$.
Let $M' = \mathcal{F}'_x$ and $M = \mathcal{F}_x$.
Then $M'$ is a $B'$-module and $M$ is a $B$-module.
Since $\mathcal{F} = i^*\mathcal{F}'$ we see that the kernel of the
surjection $M' \to M$ is $IM' = JM'$. Thus we have a short exact
sequence
$$
0 \to JM' \to M' \to M \to 0
$$
Using
Sheaves, Lemma \ref{sheaves-lemma-stalk-pullback-modules}
and
Modules, Lemma \ref{modules-lemma-stalk-tensor-product}
to identify stalks of pullbacks and tensor products we see
that the stalk at $x$ of the canonical map of the lemma is the map
$$
(J \otimes_A B) \otimes_B M = J \otimes_A M = J \otimes_{A'} M'
\longrightarrow JM'
$$
The assumption that $\mathcal{F}$ is flat over $S$ signifies that
$M$ is a flat $A$-module.

\medskip\noindent
Assume (1). Flatness implies $\text{Tor}_1^{A'}(M', A) = 0$ by
Algebra, Lemma \ref{algebra-lemma-characterize-flat}.
This means $J \otimes_{A'} M' \to M'$ is injective by
Algebra, Remark \ref{algebra-remark-Tor-ring-mod-ideal}.
Hence $J \otimes_A M \to JM'$ is an isomorphism.

\medskip\noindent
Assume (2). Then $J \otimes_{A'} M' \to M'$ is injective. Hence
$\text{Tor}_1^{A'}(M', A) = 0$ by
Algebra, Remark \ref{algebra-remark-Tor-ring-mod-ideal}.
Hence $M'$ is flat over $A'$ by
Algebra, Lemma \ref{algebra-lemma-what-does-it-mean}.
\end{proof}

\end{document}
